
\documentclass{article}
\usepackage{colortbl}
\usepackage{makecell}
\usepackage{multirow}
\usepackage{supertabular}

\begin{document}

\newcounter{utterance}

\centering \large Interaction Transcript for game `wordle\_withcritic', experiment `high\_frequency\_words\_no\_clue\_no\_critic', episode 10 with reka{-}flash{-}3.1{-}t0.0.
\vspace{24pt}

{ \footnotesize  \setcounter{utterance}{1}
\setlength{\tabcolsep}{0pt}
\begin{supertabular}{c@{$\;$}|p{.15\linewidth}@{}p{.15\linewidth}p{.15\linewidth}p{.15\linewidth}p{.15\linewidth}p{.15\linewidth}}
    \# & $\;$A & \multicolumn{4}{c}{Game Master} & $\;\:$B\\
    \hline

    \theutterance \stepcounter{utterance}  
    & & \multicolumn{4}{p{0.6\linewidth}}{
        \cellcolor[rgb]{0.9,0.9,0.9}{
            \makecell[{{p{\linewidth}}}]{
                \texttt{\tiny{[P1$\langle$GM]}}
                \texttt{You are a language wizard who likes to guess words by using the given rules.} \\
\\ 
\texttt{Welcome to Wordle! You have six attempts to guess the target word, a valid English word of five lowercase letters (a{-}z). Please use the tags "explanation:" and "guess:" to provide a concise explanation for each guess.} \\
\\ 
\texttt{For instance, if your guess is "apple", your response should be} \\
\texttt{explanation: this is a common five{-}letter English word, and I am starting my guess with this word.} \\
\texttt{guess: apple} \\
\\ 
\texttt{After each guess, your answer will be validated, and you will receive feedback indicating which letters are correct (green), which letters are correct but in the wrong position (yellow), and which letters are incorrect (red). This feedback can be useful in determining which letters to include or exclude in your next guess.} \\
\\ 
\texttt{For example, the feedback for "apple" might be:} \\
\texttt{guess\_feedback: a<yellow> p<yellow> p<green> l<yellow> e<red>} \\
\\ 
\texttt{The explanation should contain details about how the guess\_feedback is used to arrive at a new guess.} \\
\\ 
\texttt{Do not generate any other text.} \\
\\ 
\texttt{Let's begin with your first guess.} \\
\\ 
\texttt{clue: stand in (for)} \\
            }
        }
    }
    & \\ \\

    \theutterance \stepcounter{utterance}  
    & \multicolumn{4}{p{0.6\linewidth}}{
        \cellcolor[rgb]{0.9,0.9,0.9}{
            \makecell[{{p{\linewidth}}}]{
                \texttt{\tiny{[P1$\rangle$GM]}}
                \texttt{<reasoning>} \\
\texttt{Let's tackle this Wordle. The clue is "stand in (for)", which probably means the target word is a synonym for "stand in for". So, "stand in for" means to substitute or take the place of someone. The word we're looking for is a five{-}letter word that represents this action.} \\
\\ 
\texttt{First, I need to think of synonyms for "stand in for". Words like substitute, replace, take over, etc. But since it's a five{-}letter word, maybe "sub" is too short. Wait, "sub" is three letters. Let me think again. Maybe "stand in for" is being used idiomatically. Oh! "Stand in" can mean to substitute, so maybe the word is "substitute", but that's longer. Wait, the target is five letters. Let's see: substitute is 11 letters. Maybe "sub" is part of it. Wait, "sub" is three letters. Hmm.} \\
\\ 
\texttt{Alternatively, maybe the clue is a play on words. "Stand in (for)" could be a phrase where the word is formed by combining parts. Let's break it down. "Stand in" is a phrase, and "for" is another. Maybe the letters from "stand in for" form the answer. Let's see: s{-}t{-}a{-}n{-}d{-}i{-}n{-}f{-}o{-}r. But that's too many letters. Maybe the answer is a five{-}letter word that's part of that phrase. Alternatively, maybe the answer is "stand" but that's five letters. Wait, "stand" is five letters. But does "stand" relate to "stand in for"? Not exactly. Maybe "stand in for" is an idiom, and the answer is "sub" but that's three letters. Wait, maybe the answer is "sub" but the clue is "stand in (for)", so "sub" as in substitute. But "sub" is three letters. The target is five letters. Hmm.} \\
\\ 
\texttt{Wait, maybe the answer is "stand" but that's five letters. Let me check if "stand" is a valid Wordle word. Yes, it is. But does it fit the clue? "Stand in for" – maybe "stand" is part of it. Alternatively, maybe the answer is "sub" but that's three letters. Wait, the target is five letters. Let me think again. Maybe the answer is "substitute" but} \\
            }
        }
    }
    & & \\ \\

    \theutterance \stepcounter{utterance}  
    & & & \multicolumn{2}{p{0.3\linewidth}}{
        \cellcolor[rgb]{0.9,0.9,0.9}{
            \makecell[{{p{\linewidth}}}]{
                \texttt{\tiny{[GM$|$GM]}}
                \texttt{Error: The response should always start with the keyword 'explanation:'} \\
            }
        }
    }
    & & \\ \\

    \theutterance \stepcounter{utterance}  
    & & & \multicolumn{2}{p{0.3\linewidth}}{
        \cellcolor[rgb]{0.9,0.9,0.9}{
            \makecell[{{p{\linewidth}}}]{
                \texttt{\tiny{[GM$|$GM]}}
                \texttt{game\_result = ABORT} \\
            }
        }
    }
    & & \\ \\

\end{supertabular}
}

\end{document}
